\mode*
\section{Related work}
\label{related-work}

We compare our study with three lines of earlier work: debugging in
computing education, analyses of programming process data, and variation
theory and phenomenography applied to programming.

\subsection{Debugging in computing education}

The educational debugging literature up to 2008 is reviewed by
\textcite{DebuggingReview}.
On the empirical side, \textcite{DebuggingStudentPerspective} report from a
multi-institutional study that students find locating bugs harder than
fixing them, that they use a wide variety of strategies, and that the web
has emerged as the obvious first place to look for similar examples.
\Textcite{MichaeliRomeike2019} argue that debugging is a skill distinct
from general programming ability and should be taught explicitly, yet is
underrepresented both in classrooms and in computing education research.

Our study supplies what this line of work calls for but does not provide:
a \emph{theory} of what the skilled debugger does---generating the patterns
of variation---and hence of what explicit debugging instruction should
teach (\cref{debugging-as-learning,discussion}).
The finding that locating the bug is the hard part matches our analysis
that the main issue of debugging is finding the initial contrast
(\cref{debugging}); the web-first habit matches our externally-sourced
episodes, where the explanation is found rather than generated
(\cref{analysis}).

\subsection{Programming process data}

\Textcite{Jadud2006} pioneered the analysis of what he calls compilation
behaviour: observing students' repeated edit--compile cycles, with
snapshots of the program from one compilation to the next, capturing their
attempts to make programs syntactically correct.
Our edit--run cycles are the semantic sibling of Jadud's edit--compile
cycles: where compilation behaviour captures the syntactic struggle, our
variation-theoretic coding asks what the student \emph{varied} between
runs, and what that reveals about which aspects they discern.
The community has since standardised the logging of such process data in
the ProgSnap2 format \autocite{ProgSnap2}; learnlog exports its recordings
as ProgSnap2 \autocite{learnlog}, so our data and analyses remain
comparable and shareable within that ecosystem.

\subsection{Variation theory and phenomenography in programming}

\Textcite{ThuneEckerdal2009} apply phenomenography and variation theory to
engineering students' conceptions of computer programming, presenting an
outcome space of five qualitatively different ways of seeing programming,
with associated dimensions of variation.
\Textcite{BruceLearningToProgram} phenomenographically map five different
ways in which first-year students \emph{go about} learning to program,
from \enquote{following} (getting through the unit), via coding,
understanding and integrating, and problem solving, to enculturation.
Both study students' conceptions and experiences through interviews.
We differ in two ways: we apply variation theory to the recorded
\emph{acts} of debugging---the phenomenographic stance that acts express
understanding (\cref{background}), operationalised through learnlog's
edit--run data---and we connect the self-generated patterns of variation
to the approaches to learning, via \citeauthor{NCOL}'s reinterpretation
(\cref{approaches}).
Bruce et al.'s categories can be read as approach-like distinctions in the
programming context: their \enquote{following} students resemble the
surface approach, their problem solvers the deep.
Our companion paper \autocite{VTProgMisconceptions} complements these by
deriving teaching patterns from documented misconceptions.

We have not found prior work that connects the approaches to learning to
debugging specifically, nor work that looks for the patterns of variation
in programming process data.
That combination---debugging as self-directed variation, observed in
recorded edit--run cycles, and related to the deep approach---is the
contribution of this study.

\mode<presentation>{%
\begin{frame}
  \begin{block}{Related work}
    \begin{itemize}
      \item Debugging education: hard, distinct, undertaught---but no
        theory of the skill.
      \item Process data: edit--compile cycles, ProgSnap2---but analysed
        for syntax, not variation.
      \item VT/phenomenography in programming: conceptions via
        interviews---not acts, not debugging.
    \end{itemize}
  \end{block}

  \begin{idea}
    \begin{itemize}
      \item New: patterns of variation in recorded debugging, connected to
        the deep approach.
    \end{itemize}
  \end{idea}
\end{frame}%
}
